\documentclass[10pt,a4paper]{amsart}
\usepackage[margin=19mm]{geometry}
\usepackage{amsmath,amssymb,mathtools}
\usepackage[scr=rsfs,cal=euler]{mathalfa}
\usepackage{xfrac}
\usepackage[unicode,hidelinks,bookmarks=false]{hyperref}
\newcommand{\ie}{i.e.\ }
\newcommand{\cf}{cf.\ }
\newcommand{\resp}{respectively\ }
\newcommand{\Subsection}{\subsection}
\providecommand{\nobreakdash}{\nobreak\hbox{-}}
\setlength{\emergencystretch}{2em}
\multlinegap0pt
\usepackage{bm}
\usepackage{bbm}
\usepackage{dsfont}
\usepackage{xspace}
\usepackage{cancel}
\usepackage{mathtools}
\usepackage{amsthm}
\makeatletter
\@ifundefined{lemma}{\newtheorem{lemma}{Lemma}}{}
\makeatother
\newcommand{\PP}{\mathbb{P}}
\title[On the positivity of some weighted partial sums of a random ]{On the positivity of some weighted partial sums of a random multiplicative function}
\author{Marco Aymone}
\date{Source: arXiv:2408.15589v3 (CC BY 4.0)}
\begin{document}
\maketitle

\noindent\emph{Source and licence.} On the positivity of some weighted partial sums of a random multiplicative function. Version arXiv:2408.15589v3 (CC BY 4.0), distributed under the Creative Commons Attribution 4.0 International licence, as recorded in the arXiv licence metadata for this exact version and shown on the abstract page https://arxiv.org/abs/2408.15589v3. Full rights notice: licenses/arxiv-2408.15589-CC-BY-4.0.txt.\par\smallskip

\noindent\textit{Editorial scope.} Counted unit: the Lemma whose source label is \texttt{lemma estimate for the sup} in the article (source0.tex lines 316--320, source label \texttt{lemma estimate for the sup}), delivered together with the article's own complete original proof of it (source0.tex lines 321--356, one \texttt{proof} environment of 36 lines). No mathematics is written, repaired or re-derived: statement and proof are the article's own text line for line. Required context retained uncounted: lemma maximal inequality (lines 220--225); lemma explicite estimate partial sum (lines 262--266); lemma evaluacao da serie (lines 283--285). Every definition, display and label of those lines is retained unchanged. Omitted from this package and not redistributed: 1--219, 226--261, 267--282, 286--315, 357--437. Nothing inside a retained excerpt is omitted or altered: every retained body is delivered exactly as the article's lines are written, so no omission receipt of any kind accompanies this package. Citations: the retained excerpts carry no citation command at all, so no bibliography entry of the article is transcribed and no reference is redistributed. Reference adaptation: none was needed and none was made; every reference inside the delivered unit resolves inside it, so no unresolved reference or citation is produced. Printed slips kept literal: no printed word, symbol, hypothesis or number of the counted statement or of its proof was altered. Layout and class: the article's own class is \texttt{amsart} with packages including bbm, datetime, geometry, color, fancyhdr, amsmath, amssymb, amsfonts, amsthm, amstext, enumerate, amsbsy, amsopn, amscd; the wrapper is the lane's pinned \texttt{amsart} editorial wrapper at 10pt on A4, which restates the theorem-like environments the retained text needs and adds the article's own macro definitions that the retained text needs (\texttt{\textbackslash PP}), with extra wrapper packages (bm, bbm, dsfont, xspace, cancel, mathtools). The delivered numbering is the unit's own. Assets: no retained span contains a figure environment, a graphics-inclusion command, a PDF-page inclusion, a table or a TikZ picture; no image file is copied into this package, the package declares no asset dependency and no image file is redistributed with it. Licence: version arXiv:2408.15589v3 of this article is distributed under the Creative Commons Attribution 4.0 International licence, as recorded in the arXiv licence metadata for this exact version and shown on the abstract page https://arxiv.org/abs/2408.15589v3. A full rights notice is delivered at licenses/arxiv-2408.15589-CC-BY-4.0.txt. Characters: every retained excerpt is pure ASCII, so the delivered engine, XeLaTeX with its default Latin Modern text font, reports no missing character.
\begin{lemma}\label{lemma maximal inequality}
Let $f$ be a Rademacher random multiplicative function. Let $(a(n))_n$ be a sequence of non-negative real numbers such that 
$$\sum_{n=1}^\infty a^2(n)(\log n)^2<\infty.$$
Let $x\geq 1$, $m\geq 4$ and $\lambda>0$ be any fixed parameters. Then, for some absolute constant $\kappa>0$
$$\PP\left(\sup_{y>x}\left|\sum_{n>y} a(n)f(n)\right|\geq \lambda\right)\leq \frac{\kappa^{m}}{\lambda^m}\left(\sum_{n>x}\mu^2(n)a^2(n)(m-1)^{\omega(n)}\right)^{m/2}.$$
\end{lemma}

\begin{lemma}\label{lemma explicite estimate partial sum} There are positive constants $c_3>0$ and $c_5>0$ such that for all $m\geq 4$ and $x\geq 2$,
\begin{equation}\label{equation upper bound 3 omega n}
\sum_{n\leq x}\mu^2(n)(m-1)^{\omega(n)}\leq c_3mx(\log x)^{c_5m}.
\end{equation}
\end{lemma}

\begin{lemma}\label{lemma evaluacao da serie} Let $\sigma\in(1/2,1)$. Let $c_5$ be as in Lemma \ref{lemma explicite estimate partial sum}. There are absolute positive constants $c_7,c_8$ such that for all $m\geq 4$ and all $x\geq2$, 
$$\sum_{n>x}\frac{\mu^2(n)(m-1)^{\omega(n)}}{n^{2\sigma}}\leq \frac{c_7^m m^{c_5m}}{(\sigma-1/2)^{c_8m}}\frac{(\log x)^{c_5m}}{x^{2\sigma-1}}.$$
\end{lemma}

\begin{lemma}\label{lemma estimate for the sup} Let $\lambda>0$. For each $0<\theta<1$, there exists $\sigma_0>1/2$ and $\epsilon>0$ such that uniformly for all $\sigma\in(1/2,\sigma_0]$, and
$$\log x = \frac{1}{(\sigma-1/2)^{1/\theta}},$$ 
we have that for some absolute constant $\beta>0$,
$$\PP\left(\sup_{y\geq x}\left|\sum_{n>y}\frac{f(n)}{n^\sigma}\right|\geq \lambda\right)\leq \exp\left(-\beta \frac{(\log x)^{2-2\theta}}{\log\log x}+(\log\lambda ^{-1})\epsilon\frac{(\log x)^{1-\theta}}{\log\log x}\right).$$
\end{lemma}

\begin{proof}
By the previous Lemmas \ref{lemma maximal inequality} and \ref{lemma evaluacao da serie}, for any $\lambda>0$, $m\geq 4$ and all $x\geq 2$ we obtain that
$$\PP\left(\sup_{y\geq x}\left|\sum_{n>y}\frac{f(n)}{n^\sigma}\right|\geq \lambda\right)\leq \frac{\kappa^{m}}{\lambda^m}\left(\frac{c_7^m m^{c_5m}}{(\sigma-1/2)^{c_8m}}\frac{(\log x)^{c_5m}}{x^{2\sigma-1}}\right)^{m/2}.$$

Thus, we are led to optimize
\begin{multline*}
\Lambda:=\exp\bigg{(}c_9m^2\log m+c_9m^{2}\log\log x\\
+c_{10}m^2\log\left(\frac{1}{\sigma-1/2}\right)-mc_{11}(\sigma-1/2)\log x\bigg),
\end{multline*}
where $c_9,c_{10}$ and $c_{11}$ are positive constants such that our target probability is bounded above by $\lambda^{-m}\Lambda$.

In this optimization problem, we begin by choosing $x=x_\sigma$ so that 
\begin{equation}\label{equation auxiliar definition of x} 
\log x = \frac{1}{(\sigma-1/2)^{1/\theta}}, 
\end{equation} 
where $\theta\in(0,1)$ is a fixed parameter. We infer that our optimal solution occurs at
$$m= \epsilon(\sigma-1/2)\frac{\log x}{\log\log x }=\epsilon\frac{(\log x)^{1-\theta}}{\log\log x},$$
for some small fixed $\epsilon>0$ to be chosen shortly.
In particular, we have:
\begin{align*}
c_9m^2\log m &\leq c_9\epsilon^2(1-\theta)\frac{(\log x)^{2-2\theta}}{\log\log x}, \\
c_9m^2\log\log x&=c_9\epsilon^2\frac{(\log x)^{2-2\theta}}{\log\log x},\\
c_{10}m^2\log\left(\frac{1}{\sigma-1/2}\right)&=c_{10}\epsilon^2\theta\frac{(\log x)^{2-2\theta}}{\log\log x},\\
mc_{11}(\sigma-1/2)\log x&=c_{11}\epsilon\frac{(\log x)^{2-2\theta}}{\log\log x}.
\end{align*}

Collecting all these estimates, we obtain that
$$\Lambda\leq \exp\left(w(\epsilon)\frac{(\log x)^{2-2\theta}}{\log\log x}\right),$$
where 
$$w(\epsilon)=\epsilon^2( c_9(1-\theta)+c_9+c_{10}\theta)-\epsilon c_{11}:=\epsilon^2c_{12}-\epsilon c_{11}.$$

We see that $\lim_{\epsilon\to 0^+} w(\epsilon)\epsilon^{-1}=-c_{11}$, so there exists $\epsilon_0>0$ such that
$w(\epsilon_0)<0$. Thus, by setting $-\beta=w(\epsilon_0)$, our target probability is upper bounded
by
$$\frac{1}{\lambda^m}\exp\left(-\beta \frac{(\log x)^{2-2\theta}}{\log\log x}\right)=\exp\left(-\beta \frac{(\log x)^{2-2\theta}}{\log\log x}+(\log\lambda ^{-1})\epsilon\frac{(\log x)^{1-\theta}}{\log\log x} \right).$$
\end{proof}

\end{document}
